\documentclass[a4paper,11pt]{article}
\usepackage[utf8]{inputenc}
\usepackage[czech]{babel}
\usepackage{amsmath, amssymb, amsthm}
\usepackage{geometry}
\geometry{a4paper, total={170mm,257mm}, left=20mm, top=20mm}

\begin{document}

\noindent
\textbf{Křestní jméno:} \underline{\hspace{4cm}} \hfill \textbf{Příjmení:} \underline{\hspace{4cm}}

\vspace{5mm}

\noindent
\textbf{Cvičící:} \underline{\hspace{4cm}}

\vspace{10mm}

\begin{center}
    {\Large \textbf{NMAG 111 Lineární algebra 1, zimní semestr}} \\[3mm]
    {\large \textbf{1. písemná práce -- VZOROVÉ ZADÁNÍ}} \\[2mm]
    {\large (Kapitoly 2--5)}
\end{center}

\vspace{5mm}

\noindent
\textbf{Doba řešení:} 90 minut \\
\textbf{Instrukce:}
\begin{itemize}
    \item Neotvírejte dříve, než jste k tomu vyzváni dozorem!
    \item Všechny odpovědi musí být řádně zdůvodněné, není-li řečeno jinak.
    \item Žádné elektronické pomůcky včetně kalkulačky nejsou dovoleny.
    \item U úloh 1--4 pište řešení přímo do zadání, u úloh 5--7 využijte volné listy.
\end{itemize}

\vspace{5mm}

\begin{center}
\begin{tabular}{|c|c|}
\hline
\textbf{Příklad} & \textbf{Body} \\ \hline
1 [6] & \\ \hline
2 [6] & \\ \hline
3 [6] & \\ \hline
4 [6] & \\ \hline
5 [8] & \\ \hline
6 [8] & \\ \hline
7 [10] & \\ \hline
\textbf{Celkem [50]} & \\ \hline
\end{tabular}
\end{center}

\newpage

% --- PŘÍKLAD 1 ---
\noindent
\textbf{(1) [6 bodů]} Zakroužkujte správnou odpověď, nezdůvodňujte. K získání bodů za každou podúlohu je potřeba odpovědět správně na všechny tři otázky v ní. \\
\textbf{ANO} = za daných předpokladů vždy pravda, \textbf{NE} = jinak.

\vspace{3mm}

\noindent
(a) Uvažujme soustavu lineárních rovnic $Ax=b$ o $n$ neznámých nad tělesem $\mathbb{R}$, kde $A$ je čtvercová matice řádu $n$.
\begin{itemize}
    \item[\textbf{ANO / NE}] Pokud má soustava jediné řešení, pak je matice $A$ regulární.
    \item[\textbf{ANO / NE}] Pokud je $b$ nulový vektor (homogenní soustava), má soustava vždy alespoň jedno řešení.
    \item[\textbf{ANO / NE}] Pokud $rank(A) < n$, pak soustava nemá žádné řešení.
\end{itemize}

\vspace{3mm}

\noindent
(b) Nechť $V$ je vektorový prostor dimenze 3 nad tělesem $\mathbb{Z}_5$ a $u, v, w \in V$ jsou libovolné vektory.
\begin{itemize}
    \item[\textbf{ANO / NE}] Pokud jsou vektory $u, v, w$ lineárně nezávislé, pak tvoří bázi prostoru $V$.
    \item[\textbf{ANO / NE}] Vektory $u, 2u, w$ jsou vždy lineárně závislé.
    \item[\textbf{ANO / NE}] Pokud $span\{u, v\} = span\{u, w\}$, pak nutně $v = w$.
\end{itemize}

\vspace{1cm}

% --- PŘÍKLAD 2 ---
\noindent
\textbf{(2) [6 bodů]} Najděte všechna řešení soustavy lineárních rovnic nad tělesem $\mathbb{Z}_5$. Výsledek zapište ve tvaru afinního podprostoru (partikulární řešení + jádro).

$$
\left( \begin{matrix}
1 & 2 & 3 & 4 \\
2 & 4 & 1 & 3 \\
3 & 1 & 4 & 2
\end{matrix} \right)
\cdot
\left( \begin{matrix}
x_1 \\ x_2 \\ x_3 \\ x_4
\end{matrix} \right)
=
\left( \begin{matrix}
1 \\ 2 \\ 3
\end{matrix} \right)
$$

\vspace{6cm}

% --- PŘÍKLAD 3 ---
\noindent
\textbf{(3) [6 bodů]} Uvažujme matice $A$ a $B$ nad $\mathbb{R}$:
$$ A = \left( \begin{matrix} 1 & 2 \\ 3 & 6 \end{matrix} \right), \quad B = \left( \begin{matrix} 2 & 4 \\ 6 & 12 \end{matrix} \right) $$
Rozhodněte, zda existuje matice $X$ typu $2 \times 2$ taková, že $AX = B$. Pokud ano, najděte všechna taková $X$. Pokud ne, zdůvodněte proč.

\newpage

% --- PŘÍKLAD 4 ---
\noindent
\textbf{(4) [6 bodů]} V této úloze je třeba předvést podrobný výpočet se stručným komentářem zdůvodňujícím správnost postupu. \\
Najděte nad $\mathbb{Z}_7$ nějakou matici inverzní zleva k matici $A$, nebo dokažte, že taková matice neexistuje.
$$ A = \left( \begin{matrix} 1 & 5 \\ 2 & 3 \\ 4 & 1 \end{matrix} \right) $$

\vspace{8cm}

% --- PŘÍKLAD 5 ---
\noindent
\textbf{(5) [8 bodů]} Uvažujme podprostor $W$ vektorového prostoru $\mathbb{R}^4$ generovaný vektory:
$$ v_1 = (1, 1, -1, 2)^T, \quad v_2 = (2, 0, 1, 3)^T, \quad v_3 = (0, 2, -3, 1)^T $$
(a) Určete dimenzi podprostoru $W$. \\
(b) Rozhodněte, zda vektor $u = (3, 1, 0, 5)^T$ leží ve $W$. \\
(c) Doplňte (pokud je to nutné) množinu $\{v_1, v_2\}$ na bázi celého prostoru $\mathbb{R}^4$.

\newpage

% --- PŘÍKLAD 6 ---
\noindent
\textbf{(6) [8 bodů]} V této úloze máte dokázat jednodušší tvrzení ze skript. Důkaz pište pečlivě, krok za krokem. Využívejte pouze definice a základní vlastnosti operací (axiomy), případně tvrzení dokázaná v předchozích kapitolách (ta ale explicitně citujte).

\vspace{3mm}

\noindent
(a) Nechť $V$ je vektorový prostor nad tělesem $T$. Dokažte, že pro každý vektor $v \in V$ platí:
$$ (-1) \cdot v = -v $$
(kde $-v$ značí opačný vektor k vektoru $v$).

\vspace{5cm}

\noindent
(b) Nechť $A$ a $B$ jsou regulární čtvercové matice řádu $n$ nad tělesem $T$. Dokažte, že jejich součin $AB$ je také regulární matice a že platí:
$$ (AB)^{-1} = B^{-1}A^{-1} $$

\vspace{6cm}

% --- PŘÍKLAD 7 ---
\noindent
\textbf{(7) [10 bodů]} Odpovědi v následujících úlohách zdůvodněte (tj. dokažte nebo uveďte protipříklad).

\vspace{3mm}

\noindent
(a) Nechť $A$ je matice typu $m \times n$ nad $\mathbb{R}$ a $rank(A) = m$. Platí nutně, že soustava $Ax=b$ má řešení pro každou pravou stranu $b \in \mathbb{R}^m$?

\vspace{4cm}

\noindent
(b) Uvažujme vektorový prostor $V$ všech reálných polynomů stupně nejvýše 2. Tvoří množina $M = \{ p(x) \in V \mid p(0) = 1 \}$ podprostor prostoru $V$?

\vspace{4cm}

\noindent
(c) Nechť $u, v, w$ jsou lineárně nezávislé vektory ve vektorovém prostoru $V$. Jsou i vektory $u+v, v+w, u+w$ lineárně nezávislé?

\end{document}